\newpage
\section{2. ER/EER Model Design}

\subsection*{2.1 Requirements Analysis Summary}

The Smart Healthcare Platform must support:
\begin{itemize}
    \item Multiple hospitals and departments
    \item Doctors with specialties and licenses
    \item Patients with medical histories and insurance
    \item Appointments and medical records
    \item Prescriptions (weak entity dependent on appointments)
    \item Lab results and billing
    \item Security and device management
\end{itemize}

\subsection*{2.2 Entity Definitions}

\subsubsection*{1. Hospital}

\begin{tabular}{lll}
\textbf{Attribute} & \textbf{Type} & \textbf{Notes} \\
\hline
HospitalID & PK & \\
HospitalName & VARCHAR & \\
Address & VARCHAR & \\
City & VARCHAR & \\
PhoneNumber & VARCHAR & \\
\end{tabular}

\subsubsection*{2. User (Superclass for EER)}

\begin{tabular}{lll}
\textbf{Attribute} & \textbf{Type} & \textbf{Notes} \\
\hline
UserID & PK & \\
FirstName & VARCHAR & composite: Name \\
LastName & VARCHAR & composite: Name \\
Email & VARCHAR & UNIQUE \\
Phone & VARCHAR & multivalued \\
City & VARCHAR & composite: Address \\
Street & VARCHAR & composite: Address \\
PostalCode & VARCHAR & composite: Address \\
Role & VARCHAR & Patient, Doctor, Admin, SecurityAnalyst \\
\end{tabular}

\subsection*{2.3 Detailed Entity Specifications}

\subsubsection*{3. Department}

Relationship: one hospital $\rightarrow$ many departments.

\begin{tabular}{lll}
\textbf{Attribute} & \textbf{Type} & \textbf{Notes} \\
DepartmentID & PK &  \\
HospitalID & FK & references Hospital \\
DepartmentName & VARCHAR &  \\
FloorNumber & INTEGER &  \\
\end{tabular}

\subsection*{4. Patient}
\begin{tabular}{lll}
\textbf{Attribute} & \textbf{Type} & \textbf{Notes} \\
PatientID & PK &  \\
UserID & FK & references User \\
BloodType & VARCHAR &  \\
BirthDate & DATE &  \\
EmergencyContact & VARCHAR &  \\
InsuranceID & FK & references Insurance \\
\end{tabular}

\subsection*{5. Doctor}
\begin{tabular}{lll}
\textbf{Attribute} & \textbf{Type} & \textbf{Notes} \\
DoctorID & PK &  \\
UserID & FK & references User \\
DepartmentID & FK &  \\
Specialty & VARCHAR &  \\
LicenseNumber & VARCHAR & UNIQUE \\
\end{tabular}

\subsection*{6. Appointment}

This works like a junction entity.

It resolves the Doctor $\leftrightarrow$ Patient relationship.

\begin{tabular}{lll}
\textbf{Attribute} & \textbf{Type} & \textbf{Notes} \\
AppointmentID & PK &  \\
PatientID & FK &  \\
DoctorID & FK &  \\
AppointmentDate & TIMESTAMP &  \\
Status & VARCHAR &  \\
DiagnosisSummary & TEXT &  \\
\end{tabular}

\subsection*{7. Prescription (Weak Entity)}

Weak entity.

It is meaningless without an Appointment.

\begin{tabular}{lll}
\textbf{Attribute} & \textbf{Type} & \textbf{Notes} \\
PrescriptionNo & Partial Key &  \\
AppointmentID & FK & identifying relationship \\
MedicineName & VARCHAR &  \\
Dosage & VARCHAR &  \\
DurationDays & INTEGER &  \\
\end{tabular}

Composite PK:

\[(PrescriptionNo, AppointmentID)\]

\subsection*{8. LabResult}
\begin{tabular}{lll}
\textbf{Attribute} & \textbf{Type} & \textbf{Notes} \\
ResultID & PK &  \\
AppointmentID & FK &  \\
TestName & VARCHAR &  \\
ResultValue & VARCHAR &  \\
ResultDate & DATE &  \\
\end{tabular}

\subsection*{9. Insurance}
\begin{tabular}{lll}
\textbf{Attribute} & \textbf{Type} & \textbf{Notes} \\
InsuranceID & PK &  \\
ProviderName & VARCHAR &  \\
PolicyNumber & VARCHAR &  \\
CoverageType & VARCHAR &  \\
\end{tabular}

\subsection*{10. Billing}
\begin{tabular}{lll}
\textbf{Attribute} & \textbf{Type} & \textbf{Notes} \\
BillID & PK &  \\
PatientID & FK &  \\
Amount & DECIMAL &  \\
PaymentStatus & VARCHAR &  \\
BillingDate & DATE &  \\
\end{tabular}

\subsection*{11. MedicalDevice}

Cybersecurity starts here.

\begin{tabular}{lll}
\textbf{Attribute} & \textbf{Type} & \textbf{Notes} \\
DeviceID & PK &  \\
HospitalID & FK &  \\
DeviceType & VARCHAR &  \\
FirmwareVersion & VARCHAR &  \\
RiskLevel & VARCHAR &  \\
\end{tabular}

\subsection*{12. SecurityEvent}

SIEM logic.

\begin{tabular}{lll}
\textbf{Attribute} & \textbf{Type} & \textbf{Notes} \\
EventID & PK &  \\
UserID & FK &  \\
DeviceID & FK & nullable \\
EventType & VARCHAR &  \\
Severity & VARCHAR &  \\
SourceIP & VARCHAR &  \\
EventTimestamp & TIMESTAMP &  \\
\end{tabular}


Composite PK:

\[(DeviceID, VulnerabilityID)\]

\subsection*{ER Relationship Summary}

\begin{tabular}{ll}
\textbf{Relationship} & \textbf{Cardinality} \\
\hline
Hospital $\rightarrow$ Department & 1:N \\
Department $\rightarrow$ Doctor & 1:N \\
Patient $\rightarrow$ Appointment & 1:N \\
Doctor $\rightarrow$ Appointment & 1:N \\
Appointment $\rightarrow$ Prescription & 1:N \\
Appointment $\rightarrow$ LabResult & 1:N \\
Patient $\rightarrow$ Billing & 1:N \\
Hospital $\rightarrow$ MedicalDevice & 1:N \\
MedicalDevice $\leftrightarrow$ Vulnerability & N:M \\
User $\rightarrow$ SecurityEvent & 1:N \\
\end{tabular}

\subsection*{EER Model: Enhanced Features}

\subsection*{Overview}
Now the EER part.

\subsection*{Specialization / Generalization}
Superclass

User

Subclasses
Patient
Doctor
Admin
SecurityAnalyst

This structure shows:

inheritance
specialization
ISA hierarchy

\subsection*{EER Features}

EER Özellikleri
Type:

Disjoint Specialization

A user cannot be both:

the Patient
the Doctor

at the same time.

Participation:

Total Participation

Every User must have a role.

\subsection*{Composite Attribute Example}

User:

Name:

FirstName
LastName

Address:

City
Street
PostalCode
\subsection*{Multivalued Attribute}

Patient:

PhoneNumbers

During relational mapping, this is separated into its own table:

PatientPhone

\begin{tabular}{ll}
PatientID & PhoneNumber \\
\end{tabular}

\subsection*{Weak Entity}

Prescription

Owner entity:

Appointment

Partial key:

PrescriptionNo
\subsection*{Category / Union Type (Optional Bonus)}
StaffAccess

Union:

Doctor
Admin
SecurityAnalyst

It is optional; do not overcomplicate it.

\section{Diagram Checklist}
ER diagram should show:
\begin{itemize}
    \item PK
    \item FK
    \item cardinality
    \item weak entity
    \item total participation
\end{itemize}

EER diagram should additionally show:
\begin{itemize}
    \item ISA hierarchy
    \item specialization
    \item inheritance
    \item disjoint/overlap
    \item total/partial participation
\end{itemize}

This structure is enough for the ER/EER part of the assignment.

% (Removed duplicate positioning-based TikZ figures — coordinate-based figures are above)

\newpage
\section{3. Relational Schema Mapping}

Convert the ER/EER design into normalized relational tables:

\subsection*{3.1 Relational Schema}

\textbf{User} (UserID, FirstName, LastName, Email, Phone, City, Street, PostalCode, Role)
\begin{itemize}
    \item PK: UserID
    \item Composite attribute: Name (FirstName, LastName), Address (City, Street, PostalCode)
    \item Multivalued Phone handled in UserPhone table
\end{itemize}

\textbf{Hospital} (HospitalID, HospitalName, Address, City, PhoneNumber)
\begin{itemize}
    \item PK: HospitalID
\end{itemize}

\textbf{Department} (DepartmentID, HospitalID, DepartmentName, FloorNumber)
\begin{itemize}
    \item PK: DepartmentID
    \item FK: HospitalID $\rightarrow$ Hospital
\end{itemize}

\textbf{Doctor} (DoctorID, UserID, DepartmentID, Specialty, LicenseNumber)
\begin{itemize}
    \item PK: DoctorID
    \item FK: UserID $\rightarrow$ User
    \item FK: DepartmentID $\rightarrow$ Department
    \item UNIQUE: LicenseNumber
\end{itemize}

\textbf{Patient} (PatientID, UserID, BloodType, BirthDate, EmergencyContact, InsuranceID)
\begin{itemize}
    \item PK: PatientID
    \item FK: UserID $\rightarrow$ User
    \item FK: InsuranceID $\rightarrow$ Insurance
\end{itemize}

\textbf{Appointment} (AppointmentID, PatientID, DoctorID, AppointmentDate, Status, DiagnosisSummary)
\begin{itemize}
    \item PK: AppointmentID
    \item FK: PatientID $\rightarrow$ Patient
    \item FK: DoctorID $\rightarrow$ Doctor
\end{itemize}

\textbf{Prescription} (PrescriptionNo, AppointmentID, MedicineName, Dosage, DurationDays)
\begin{itemize}
    \item PK: (PrescriptionNo, AppointmentID) — Composite Key
    \item FK: AppointmentID $\rightarrow$ Appointment
    \item Weak entity dependent on Appointment
\end{itemize}

\textbf{LabResult} (ResultID, AppointmentID, TestName, ResultValue, ResultDate)
\begin{itemize}
    \item PK: ResultID
    \item FK: AppointmentID $\rightarrow$ Appointment
\end{itemize}

\textbf{Insurance} (InsuranceID, ProviderName, PolicyNumber, CoverageType)
\begin{itemize}
    \item PK: InsuranceID
\end{itemize}

\textbf{Billing} (BillID, PatientID, Amount, PaymentStatus, BillingDate)
\begin{itemize}
    \item PK: BillID
    \item FK: PatientID $\rightarrow$ Patient
\end{itemize}

\textbf{MedicalDevice} (DeviceID, HospitalID, DeviceType, FirmwareVersion, RiskLevel)
\begin{itemize}
    \item PK: DeviceID
    \item FK: HospitalID $\rightarrow$ Hospital
\end{itemize}

\textbf{SecurityEvent} (EventID, UserID, DeviceID, EventType, Severity, SourceIP, EventTimestamp)
\begin{itemize}
    \item PK: EventID
    \item FK: UserID $\rightarrow$ User
    \item FK: DeviceID $\rightarrow$ MedicalDevice (nullable)
\end{itemize}

\textbf{UserPhone} (UserID, PhoneNumber)
\begin{itemize}
    \item PK: (UserID, PhoneNumber)
    \item FK: UserID $\rightarrow$ User
\end{itemize}

\newpage
\section{4. Functional Dependencies}

\subsection*{4.1 Hospital Table}

\begin{itemize}
    \item HospitalID $\rightarrow$ HospitalName, Address, City, PhoneNumber
    \item All attributes depend solely on HospitalID
\end{itemize}

\subsection*{4.2 Doctor Table}

\begin{itemize}
    \item DoctorID $\rightarrow$ UserID, DepartmentID, Specialty, LicenseNumber
    \item UserID $\rightarrow$ FirstName, LastName, Email (via User table)
    \item DepartmentID $\rightarrow$ DepartmentName, FloorNumber (via Department table)
    \item LicenseNumber $\rightarrow$ DoctorID (candidate key)
\end{itemize}

\subsection*{4.3 Patient Table}

\begin{itemize}
    \item PatientID $\rightarrow$ UserID, BloodType, BirthDate, EmergencyContact, InsuranceID
    \item UserID $\rightarrow$ FirstName, LastName, Email (via User table)
    \item InsuranceID $\rightarrow$ ProviderName, PolicyNumber, CoverageType (via Insurance table)
\end{itemize}

\subsection*{4.4 Appointment Table}

\begin{itemize}
    \item AppointmentID $\rightarrow$ PatientID, DoctorID, AppointmentDate, Status, DiagnosisSummary
    \item PatientID $\rightarrow$ BloodType, BirthDate (via Patient table)
    \item DoctorID $\rightarrow$ Specialty, LicenseNumber (via Doctor table)
\end{itemize}

\subsection*{4.5 Prescription Table}

\begin{itemize}
    \item (PrescriptionNo, AppointmentID) $\rightarrow$ MedicineName, Dosage, DurationDays
    \item AppointmentID $\rightarrow$ PatientID, DoctorID (via Appointment table)
\end{itemize}

\newpage
\section{5. Normalization Analysis}

\subsection*{5.1 First Normal Form (1NF)}

\textbf{Status:} \checkmark~Already in 1NF

All tables have atomic values (no composite or multivalued attributes in main tables).
\begin{itemize}
    \item User: Phone moved to UserPhone table
    \item Address attributes (City, Street, PostalCode) kept as separate columns (acceptable for 1NF)
\end{itemize}

\subsection*{5.2 Second Normal Form (2NF)}

\textbf{Status:} \checkmark~Already in 2NF

All non-key attributes are fully dependent on the primary key.

\begin{itemize}
    \item Doctor table: DoctorID is PK; Specialty and LicenseNumber depend fully on DoctorID
    \item Patient table: PatientID is PK; BloodType, BirthDate depend fully on PatientID
    \item Prescription table: Composite PK (PrescriptionNo, AppointmentID); MedicineName, Dosage depend on the full key
\end{itemize}

\subsection*{5.3 Third Normal Form (3NF)}

\textbf{Status:} \checkmark~Already in 3NF

No transitive dependencies. Non-key attributes do not depend on other non-key attributes.

\begin{itemize}
    \item Doctor: Specialty depends directly on DoctorID, not on another non-key attribute
    \item Patient: BloodType depends directly on PatientID
    \item Appointment: Status depends on AppointmentID, not transitively through PatientID
\end{itemize}

\subsection*{5.4 Boyce-Codd Normal Form (BCNF)}

\textbf{Status:} \checkmark~BCNF-compliant for critical tables

Every determinant is a candidate key.

\begin{itemize}
    \item Doctor: LicenseNumber is a candidate key; all FDs originate from candidate keys
    \item Hospital: HospitalID is the only determinant and is the PK
    \item Insurance: InsuranceID is the only determinant and is the PK
\end{itemize}

\subsection*{5.5 Normalization Summary Table}

\begin{tabular}{lcccc}
\textbf{Table} & \textbf{1NF} & \textbf{2NF} & \textbf{3NF} & \textbf{BCNF} \\
\hline
Hospital & \checkmark & \checkmark & \checkmark & \checkmark \\
Doctor & \checkmark & \checkmark & \checkmark & \checkmark \\
Patient & \checkmark & \checkmark & \checkmark & \checkmark \\
Appointment & \checkmark & \checkmark & \checkmark & \checkmark \\
Prescription & \checkmark & \checkmark & \checkmark & \checkmark \\
Department & \checkmark & \checkmark & \checkmark & \checkmark \\
\end{tabular}

All tables are normalized up to BCNF, ensuring data integrity and minimizing redundancy.
